% !TEX TS-program = xelatex
% !TEX encoding = UTF-8 Unicode
\documentclass[10pt,aspectratio=1610]{beamer}
%\documentclass[10pt]{beamer}

\usetheme[progressbar=frametitle]{metropolis}
\usepackage{appendixnumberbeamer}


\usepackage{booktabs}
\usepackage[scale=2]{ccicons}

\usepackage{pgfplots}
\usepgfplotslibrary{dateplot}

\usepackage{xspace}

\usepackage{bm}

% Color definitions 
%\usepackage[dvipsnames]{xcolor}

% Remove "Figure #" from figures
\setbeamertemplate{caption}{\raggedright\insertcaption\par}

% Custom commands
% Metropolis theme name 
\newcommand{\themename}{\textbf{\textsc{metropolis}}\xspace}

% Derivatives
\newcommand{\pfrac}[3][]{\frac{\partial^{#1} #2}{\partial #3^{#1}}}
\newcommand{\dd}[3][]{\frac{\mathrm{d}^{#1} #2}{\mathrm{d} #3^{#1}}}
\newcommand{\pp}[2][]{\partial^{#1}_{#2}}


\newcounter{example}
\stepcounter{example}
\newcounter{exercise}
\stepcounter{exercise}

\renewcommand{\vec}{\mathbf}
\newcommand{\varone}{\vec{x}}
\newcommand{\argone}{(t,\varone)}
\newcommand{\R}{\mathbb{R}}
\newcommand{\CC}{\mathbb{C}}
\newcommand{\Q}{\mathbb{Q}}
\newcommand{\Z}{\mathbb{Z}}
\newcommand{\N}{\mathbb{N}}
\newcommand{\norm}[1]{\left\lVert#1\right\rVert}
\newcommand{\fone}{\vec{f}}
\newcommand{\ftwo}{\vec{g}}
\newcommand{\fthr}{\vec{h}}
\newcommand{\ffor}{\bm{\phi}}
\newcommand{\opone}{\mathcal{L}}
\newcommand{\optwo}{\mathcal{G}}
\renewcommand{\Re}{\operatorname{Re}}
\newcommand{\iprod}[2]{\langle #1,#2 \rangle}
\newcommand{\diff}{\mathrm{d}}
\newcommand{\fouriert}{\mathcal{F}}

\title{Heat and wave equations}
\subtitle{Introduction in 1d}
\date{16/2/2021}
\date{}
%\author{18.303 Linear Partial Differential Equations: Analysis and Numerics}
\institute{18.303 Linear Partial Differential Equations: Analysis and Numerics}
\titlegraphic{\hfill\includegraphics[height=2em]{../MIT-logo.pdf}}

\begin{document}
	
	\maketitle
	
%	\begin{frame}{Table of contents}
%		\setbeamertemplate{section in toc}[sections numbered]
%		\tableofcontents%[hideallsubsections]
%	\end{frame}

%%%%%%%%%%%%%%%
%\section{Finite differences}
%%%%%%%%%%%%%%%

%--------------------
% Adjoints 

\begin{frame}{Operator adjoints, redux}
	We recall that the adjoint of an operator $ \opone $, $ \opone^{*} $, is defined using the inner product as 
	\[ \iprod{\fone}{\opone \ftwo}  = \iprod{\opone^{*}\fone}{\ftwo}\]
	for all $ \fone, \ftwo \in V $. If $ \opone = \opone^{*} $ we say the operator is self-adjoint. 
\end{frame}

\begin{frame}{Operator adjoints, redux}
	\textbf{Example \examplen:} Let us again think about functions $ u: [0,1] \to \R $ with Dirichlet boundaries ($ u(0)=u(1)=0 $). We will look at the Laplace operator $ \Delta =\dd[2]{}{x}$. Writing the inner product for some $ f $ and $ g $ gives
	\[ \iprod{f}{\Delta g} = \int_{0}^{1} f(x) \dd[2]{}{x} g(x) \diff x. \]
	
	\pause
	We can use integration by part giving 
	\[ \left. f(x) \dd{}{x} g(x) \right|_{0}^{1} -  \int_{0}^{1} \dd{}{x} f(x) \dd{}{x} g(x) \diff x. \]
	
	\pause
	We notice that the first term is zero due to the boundary conditions. Repeating integration by parts gives
	\[ \iprod{f}{\Delta g} = \int_{0}^{1} \dd[2]{}{x} f(x)  g(x) \diff x = \iprod{\Delta f}{g} \]
	i.e. the Laplacian is self-adjoint with these boundaries.
\end{frame}

\begin{frame}{Positive and negative definite operators}
	\begin{mydefinition}
		The self-adjoint operator $ \opone $ is \alert{positive definite} iff $ \iprod{\fone}{\opone \fone} > 0 $ for all $ 0 \neq \fone \in V $.
	\end{mydefinition}
	If the equality is not strict (i.e. $ \geq $), the operator is called \alert{positive semidefinite}. 
	
	\pause
	If the inner product is negative we have the corresponding concepts \alert{negative definite} and \alert{negative semidefinite} operators.
	
	\pause
	\textbf{Example \examplen :} The Laplacian for the vectors $ u: [0,1] \to \R $ with Dirichlet boundaries is negative definite. 
	
	\pause
	\textbf{Proof:} 
	\[ \iprod{u}{\Delta u} = \int_{0}^{1} u(x) \dd[2]{}{x} u(x) \diff x = -\int_{0}^{1} \dd{u(x)}{x} \dd{u(x)}{x} = -\norm{\dd{u}{x}}^2 < 0. \]
	
	Here we used integration by part as before. Why can't the norm be zero?
\end{frame}

\begin{frame}{Superposition principle}
	Assume $ \fone $ and $ \ftwo $ are in the null space of the operator $ \opone $ i.e. ($ \opone \fone = \opone \ftwo = 0 $). 
	
	\pause
	This implies that for any constants $ \alpha $ and $ \beta  $, $ \alpha \fone + \beta \ftwo  $ is also in the null space. 
	
	\pause
	Often linear differential equations are in fact problems for finding the null space of the operator. 
	
	\pause
	Boundary conditions give the necessary information for fixing all the coefficients in front of the individual solutions for the differential equation.
\end{frame}

\begin{frame}{Heat equation}
	We look at functions $ u: \R_{+} \times [0,1] \to \R $. These are functions $ u(t,x) $ defined at non-negative times and on the interval $x \in  [0,1] $. Let us assume Dirichlet boundaries again.
	
	\pause
	\metroset{block=fill}
	\begin{block}{\centering Heat equation with Dirichlet boundaries}
		\begin{align*}
			&\pfrac{u(t,x)}{t} = \pfrac[2]{u(t,x)}{x}; \\
			&u(t,0) = u(t,1) = 0,\; t \in \R_{+}; \\
			&u(0,x) = u_0(x), \; x \in [0,1].
		\end{align*}
	\end{block}
	
	\pause
	\textbf{Important!} When the domain of the function is a Cartesian product of two spaces and the boundary conditions are defined \emph{independently} for all variables, a linear differential equation is \emph{separable}. This means that the individual solutions are of the form $ a_n (t) \phi_n (x) $.
\end{frame}

\begin{frame}
	Plugging in the separated ansatz gives 
	\[ \pfrac{a_n(t)}{t} \phi_n (x) = a_n(t) \pfrac[2]{\phi_n(x)}{x}. \]
	
	\pause
	Dividing by the $ a_n $ (let's not worry at this point if this can be zero) gives
	\[ \frac{\dot{a}_{n}(t)}{a_{n}(t)} = \frac{\phi_n''(x)}{\phi_n(x)}. \]
	
	\pause
	We note that we can adjust $ t $ and $ x $ \emph{independently}. The only possible way this equation holds for all $ t $ and $ x $ is if both sides are constant. We get two eigenvalue problems:
	
	\pause
	\begin{align*}
		\phi_n''(x) &= -\lambda_n^2 \phi_n,  \\
		\dot{a}_{n}(t) &= -\lambda_n^2 a_n(t).
	\end{align*}
	
	We know that the solution to the first one is \alert{$ \phi_n(x) = \sin(n \pi x) $} (at any given time we have Dirichlet boundaries). This also gives \alert{$ \lambda_n = n \pi $}.
\end{frame}

\begin{frame}
	The second equation is a simple ordinary differential equation, which we can solve by writing it as 
	\[ \dd{\log(a_n(t))}{t} = - \lambda_n^2. \]
	
	\pause
	We have the initial data $ u(0,x) = u_0(x) $ suggesting we integrate the equation above from $ 0 $ to $ t $. This gives
	
	\pause
	\[ \begin{split}
		\int_{0}^{t}  \dd{\log(a_n(t'))}{t'} \diff t' &= \int_{0}^{t} -\lambda_n \diff t, \\
		\log(a_n(t))-\log(a_n(0)) &= -\lambda_n^2 t, \\
		\log(a_n(t)/a_n(0)) &= -\lambda_n^2 t.
	\end{split} \]
	
	\pause
	Solving for $ a_n(t) $ gives
	\begin{equation}\label{heat_coefficients}
		a_n(t) = a_n(0) e^{-\lambda_n^2 t}.
	\end{equation}
	
\end{frame}

\begin{frame}
	Now we have the general solution (according to the superposition principle) given as
	\metroset{block=fill}
	\begin{block}{\centering Heat equation with Dirichlet boundaries, solution}
		\begin{align*}
			u(t,x) &= \sum_{n=1}^{\infty} a_n(t) \phi_n(x), \; \text{where} \\
			a_n(t) &= a_n(0) e^{-\lambda_n^2 t}, \; \phi_n(x) = \sin(n \pi x).
		\end{align*}
	\end{block}

	\pause
	How do we get the coefficients $ a_n(0) $?
	
	\pause
	At time $ t=0 $ we have 
	\[ u(0,x) = u_0 (x) = \sum_{n=1}^{\infty} a_n(0) \sin(n \pi x). \]
	
	\pause
	What we do is we express $ u_0  $ in the basis $ \{ \phi_n \}_{n=1}^\infty $ giving the coefficients $ a_n(0) $. If $ u_0 = \sum_{n=1}^{\infty} \hat{u}_0^{(n)} \phi_n $, then $ a_n(0) =\hat{u}_0^{(n)} $.
\end{frame}

\begin{frame}{Exercise \exercisen}
	Solve the heat equation on the same domain with the same initial condition and the boundary condition 
	\[ u(t,0) = 0, \; u'(t,1) = 0. \]
	
	How is the solution different?
\end{frame}

\begin{frame}{Wave equation}
	Again, $ u: \R_{+} \times [0,1] \to \R $ and we assume Dirichlet boundaries.
	
	\pause
	\metroset{block=fill}
	\begin{block}{\centering Wave equation with Dirichlet boundaries}
		\begin{align*}
			&\pfrac[2]{u(t,x)}{t} = \pfrac[2]{u(t,x)}{x}; \\
			&u(t,0) = u(t,1) = 0,\; t \in \R_{+}; \\
			&u(0,x) = u_0(x),\; \dot{u}(0,x) = v_0 (x), \; x \in [0,1].
		\end{align*}
	\end{block}

	Notice that we need more initial data since the order of the PDE in time is 2. 

%	\textbf{Important!} When the domain of the function is a Cartesian product of two spaces and the boundary conditions are defined \emph{independently} for all variables, a linear differential equation is \emph{separable}. This means that the individual solutions are of the form $ a_n (t) \phi_n (x) $.
\end{frame}

\begin{frame}
	Again, the system can be separated and we get the eigenvalue problems
	\begin{align*}
		\phi_n''(x) &= -\lambda_n^2 \phi_n,  \\
		\ddot{a}_{n}(t) &= -\lambda_n^2 a_n(t).
	\end{align*}

	\pause
	The general single (real solution) for the latter equation is 
	\[ a_n(t) = \alpha_n \cos(\lambda_n t) + \beta_n \sin(\lambda_n t). \]
\end{frame}

\begin{frame}
	\metroset{block=fill}
	\begin{block}{\centering Wave equation with Dirichlet boundaries, solution}
		\begin{align*}
			u(t,x) &= \sum_{n=1}^{\infty} a_n(t) \phi_n(x), \; \text{where} \\
			a_n(t) &= \alpha_n \cos(\lambda_n t) + \beta_n \sin(\lambda_n t), \; \phi_n(x) = \sin(n \pi x)(=\sin(\lambda_n x)).
		\end{align*}
	\end{block}

	\pause
	We get the coefficients $ \alpha_n $ and $ \beta_n $ from the initial condition:
	\begin{align*}
		u_0 (x) &= \sum_{n=1}^{\infty} \alpha_n \phi_n(x), \\
		v_0 (x) &= \dot{u}(0,x) = \sum_{n=1}^{\infty} \beta_n \lambda_n \phi_n(x).
	\end{align*}

	Let $ \hat{u}_0^{(n)} $ and $ \hat{v}_0^{(n)} $ be the sine series coefficients of $ u $ and $ v $, respectively. We have 
	\[ \alpha_n = \hat{u}_0^{(n)} ,\; \beta_n = \hat{v}_0^{(n)}/\lambda_n. \]
\end{frame}

\begin{frame}{Some observations}
	\begin{itemize}[<+->]
		\item Notice how solving the time dependent part of the equation reduced to solving an ordinary (simple) differential equation
		\item The same methodology carries over to higher dimensional complicated structures
		\item Often we have to solve the eigenvalue of the Laplacian in a complicated domain (on a computer) but after this, solving for the time evolution of the spatial eigenmodes (here $ \phi_n $) is pretty much trivial
		\item This idea can be generalized to other dynamical systems: this property follows from the fact that time is an \emph{independent} parameter and doesn't mix with the spatial part of the boundary value problem	
	\end{itemize}
\end{frame}

\begin{frame}{Notation}
	\begin{table}
		\caption{Some notation for differential operators}
		\begin{tabular}{lccc}
			\toprule
			Name & Notation & Alternative notation & Cartesian coordinates \\
			\midrule
			Total differential & $ \diff f/\diff x $ & $ f'(x) $ (1d) & - \\
			Partial differential & $ \partial f/\partial x $ & $ \pp{x} f $, $ f_x $ & -\\
			Time derivative & $ \diff f/\diff t $ & $ \dot{f} $ & - \\
			Gradient & $ \nabla \cdot f $ & $ \operatorname{grad}(f) $ & $ (\pp{x}f,\pp{y}f) $ \\
			Divergence & $ \nabla \cdot \fone $ & $ \operatorname{div}(\fone) $ & $ \pp{x} f_x + \pp{y} f_y $ \\
			Curl (2d) & $ \nabla \times \fone $ & $ \nabla \wedge \fone $ & $ \pp{x}f_y - \pp{y} f_x $ \\
			Laplacian & $ \Delta f$ & $ \nabla^2 f $ &  $ (\pp[2]{x} + \pp[2]{y})f $ \\
			\bottomrule		
		\end{tabular}
	\end{table}
\end{frame}

\end{document}
